% !TEX root = ../report.tex

\ssr{ЗАКЛЮЧЕНИЕ}

В ходе выполнения курсового проекта была разработана база данных онлайн-маркетплейса и приложение доступа к ней.

В аналитической части рассмотрена предметная область, выполнено сравнение существующих решений, сформулированы требования и ограничения проектируемой системы. Выделены роли пользователей, сценарии взаимодействия, сущности и связи между ними. Для хранения данных выбрана реляционная модель.

В конструкторской части спроектированы таблицы базы данных, ограничения целостности, индексы, ролевая модель доступа, триггеры и хранимая функция. Выполнена проверка соответствия схемы третьей нормальной форме.

В технологической части описана реализация на PostgreSQL, Go и Redis. Реализованы миграции схемы, роли СУБД, триггеры, хранимая функция, REST API, серверный Web-интерфейс и кэширование часто читаемых данных по схеме Cache-Aside.

В исследовательской части выполнены измерения времени выполнения SQL-запросов при разных наборах индексов и времени HTTP-ответа при отключённом и прогретом Redis-кэше. Исследование показало, что индекс ускоряет запрос только при соответствии структуре условия и сортировки, а Redis-кэш снижает время ответа за счёт сокращения обращений к PostgreSQL.

Поставленная цель достигнута. Разработана база данных маркетплейса и приложение доступа к ней. Исследованы выбранные способы ускорения чтения данных.

\clearpage
